\begin{problem}[Same topological point, different infinitesimal thickness]
For $n\ge1$, let
\[
A_n=k[t]/(t^n).
\]
Prove that all $\Spec A_n$ are homeomorphic one-point spaces.  Show that $A_m\not\cong A_n$
when $m\ne n$ by using nilpotency order, and compute $(A_n)_{\mathrm{red}}$.
\end{problem}
\paragraph{Why this problem matters.}
It produces an infinite family of distinct affine structures with exactly the same point-set
topology.
\paragraph{Useful ideas before starting.}
Every prime contains $\bar t$ because it is nilpotent.  The quotient by $(\bar t)$ is $k$.
\paragraph{Guided hints.}
Use the least positive power of the nilradical that becomes zero as an isomorphism invariant.
\begin{solution}
Every prime of $A_n$ contains $\bar t$.  Since $(\bar t)$ is maximal and
$A_n/(\bar t)\cong k$, it is the unique prime.  Thus each spectrum is one point.  The
nilradical is $(\bar t)$ and the reduction is $k$.
For $n\ge2$, the nilradical satisfies
\[
(\bar t)^n=0,\qquad (\bar t)^{n-1}\ne0.
\]
This nilpotency index is preserved by ring isomorphism, so $A_m\not\cong A_n$ for $m\ne n$.
\end{solution}
\paragraph{What happened mathematically.}
Topology records support; nilpotency order records thickness.
\paragraph{Extensions.}
Determine all quotient maps $A_n\twoheadrightarrow A_m$ that send $\bar t$ to $\bar t$.
\paragraph{Provenance.}
New synthesis problem built from the legacy doubled-point comparison in
\emph{Theory of Commutative Algebra 13}.
